% GENERATED FILE. Edit canonical lessons, then run npm run generate:books.
% canonical-source-hash: fnv1a64:b5a2e1542c843806
% canonical-lessons: FR-C331-moule, FR-C331-enfourner, FR-C331-cuisson, FR-C331-dorer, FR-C331-gratiner, FR-R331-revision-324-331

\chapter{Au four}
\label{ch:fr-au-four}
\hlchaptermodality{331}

\begin{chapteropening}
\textbf{By the end of this chapter:} \emph{I can describe baking: the tin, putting a dish in the oven, the cooking time, turning golden and browning on top.}

Complete the last lesson of chapter 331: I can describe baking: the tin, putting a dish in the oven, the cooking time, turning golden and browning on top.
\end{chapteropening}

\section[le moule]{le moule — the baking tin, the mould}

\label{lesson:FR-C331-moule}



\begin{quote}
Before the new one: say the French for to simmer, then the French for the casserole dish.
\end{quote}

\subsection*{You'll want to know: le moule}
\textbf{le moule} — \textquotedblleft{}the baking tin, the mould\textquotedblright{}.

\begin{quote}
Elle verse la pâte dans le moule.
\end{quote}

\begin{quote}
\emph{She pours the mixture into the tin.}
\end{quote}

Masculine; \emph{la moule} is a mussel.

\begin{cousinweb}
\emph{moule} comes from Old French \emph{modle}, inherited from Latin \emph{modulus} \textquotedblleft{}a small measure\textquotedblright{} (EtymDB), built on \emph{modus} \textquotedblleft{}measure\textquotedblright{}.
\end{cousinweb}

\begin{grammarlens}[title={What you've built}]
One more word for telling, explaining and describing at B1.
\end{grammarlens}

\subsection*{Your turn}
\begin{itemize}
\raggedright
  \item \emph{Say it:} \emph{le moule}
  \item \emph{Say it:} \emph{le moule}, once more
  \item \emph{Say it:} a sentence of your own with \emph{le moule}
  \item \emph{Recall:} say the French for to simmer, then the French for the casserole dish, then say \emph{le moule} again
\end{itemize}

\begin{tcolorbox}[breakable,colback=teal!4,colframe=teal!35!black,title={Before you move on}]
What does le moule mean? (\textquotedblleft{}The baking tin, the mould\textquotedblright{}.) Say it once more.
\end{tcolorbox}
\section[enfourner]{enfourner — to put in the oven}

\label{lesson:FR-C331-enfourner}



\begin{quote}
Before the new one: say the French for the casserole dish, then the French for the baking tin.
\end{quote}

\subsection*{You'll want to know: enfourner}
\textbf{enfourner} — \textquotedblleft{}to put in the oven\textquotedblright{}.

\begin{quote}
Il enfourne le gâteau pour quarante minutes.
\end{quote}

\begin{quote}
\emph{He puts the cake in the oven for forty minutes.}
\end{quote}

Informal: to gobble up food.

\begin{cousinweb}
\emph{enfourner} is built on \emph{en-} and \emph{four} \textquotedblleft{}oven\textquotedblright{}, which is inherited from Latin \emph{furnus} \textquotedblleft{}an oven\textquotedblright{} (Lewis \& Short).
\end{cousinweb}

\begin{grammarlens}[title={What you've built}]
One more word for telling, explaining and describing at B1.
\end{grammarlens}

\subsection*{Your turn}
\begin{itemize}
\raggedright
  \item \emph{Say it:} \emph{enfourner}
  \item \emph{Say it:} \emph{enfourner}, once more
  \item \emph{Say it:} a sentence of your own with \emph{enfourner}
  \item \emph{Recall:} say the French for the casserole dish, then the French for the baking tin, then say \emph{enfourner} again
\end{itemize}

\begin{tcolorbox}[breakable,colback=teal!4,colframe=teal!35!black,title={Before you move on}]
What does enfourner mean? (\textquotedblleft{}To put in the oven\textquotedblright{}.) Say it once more.
\end{tcolorbox}
\section[la cuisson]{la cuisson — the cooking, the cooking time}

\label{lesson:FR-C331-cuisson}



\begin{quote}
Before the new one: say the French for the baking tin, then the French for to put in the oven.
\end{quote}

\subsection*{You'll want to know: la cuisson}
\textbf{la cuisson} — \textquotedblleft{}the cooking, the cooking time\textquotedblright{}.

\begin{quote}
La cuisson du gâteau dure quarante minutes.
\end{quote}

\begin{quote}
\emph{The cake takes forty minutes to cook.}
\end{quote}

For meat: \emph{quelle cuisson ?} — how would you like it cooked?

\begin{cousinweb}
\emph{cuisson} is usually traced to Latin \emph{coctiō} \textquotedblleft{}a cooking\textquotedblright{}, from \emph{coquere} \textquotedblleft{}to cook\textquotedblright{} (Lewis \& Short), reshaped after \emph{cuire}.
\end{cousinweb}

\begin{grammarlens}[title={What you've built}]
One more word for telling, explaining and describing at B1.
\end{grammarlens}

\subsection*{Your turn}
\begin{itemize}
\raggedright
  \item \emph{Say it:} \emph{la cuisson}
  \item \emph{Say it:} \emph{la cuisson}, once more
  \item \emph{Say it:} a sentence of your own with \emph{la cuisson}
  \item \emph{Recall:} say the French for the baking tin, then the French for to put in the oven, then say \emph{la cuisson} again
\end{itemize}

\begin{tcolorbox}[breakable,colback=teal!4,colframe=teal!35!black,title={Before you move on}]
What does la cuisson mean? (\textquotedblleft{}The cooking, the cooking time\textquotedblright{}.) Say it once more.
\end{tcolorbox}
\section[dorer]{dorer — to turn golden, to brown; to gild}

\label{lesson:FR-C331-dorer}



\begin{quote}
Before the new one: say the French for to put in the oven, then the French for the cooking.
\end{quote}

\subsection*{You'll want to know: dorer}
\textbf{dorer} — \textquotedblleft{}to turn golden, to brown; to gild\textquotedblright{}.

\begin{quote}
Le poulet dore lentement au four.
\end{quote}

\begin{quote}
\emph{The chicken is slowly turning golden in the oven.}
\end{quote}

\emph{Faire dorer les oignons} — to brown the onions.

\begin{cousinweb}
\emph{dorer} is inherited through Old French \emph{dorer} (EtymDB) from Latin \emph{deaurāre} \textquotedblleft{}to gild\textquotedblright{}, built on \emph{aurum} \textquotedblleft{}gold\textquotedblright{} (Lewis \& Short).
\end{cousinweb}

\begin{grammarlens}[title={What you've built}]
One more word for telling, explaining and describing at B1.
\end{grammarlens}

\subsection*{Your turn}
\begin{itemize}
\raggedright
  \item \emph{Say it:} \emph{dorer}
  \item \emph{Say it:} \emph{dorer}, once more
  \item \emph{Say it:} a sentence of your own with \emph{dorer}
  \item \emph{Recall:} say the French for to put in the oven, then the French for the cooking, then say \emph{dorer} again
\end{itemize}

\begin{tcolorbox}[breakable,colback=teal!4,colframe=teal!35!black,title={Before you move on}]
What does dorer mean? (\textquotedblleft{}To turn golden, to brown; to gild\textquotedblright{}.) Say it once more.
\end{tcolorbox}
\section[gratiner]{gratiner — to brown on top, to cook with a crust}

\label{lesson:FR-C331-gratiner}



\begin{quote}
Before the new one: say the French for the cooking, then the French for to turn golden.
\end{quote}

\subsection*{You'll want to know: gratiner}
\textbf{gratiner} — \textquotedblleft{}to brown on top, to cook with a crust\textquotedblright{}.

\begin{quote}
Elle fait gratiner les pâtes avec du fromage.
\end{quote}

\begin{quote}
\emph{She browns the pasta with cheese on top.}
\end{quote}

\emph{Un gratin} — a dish browned in the oven, often with cheese.

\begin{cousinweb}
\emph{gratiner} is built on \emph{gratin}, the crust scraped from the dish, from the verb meaning \textquotedblleft{}to scrape\textquotedblright{}, which is usually traced to a Frankish word (EtymDB lists no parent).
\end{cousinweb}

\begin{grammarlens}[title={What you've built}]
One more word for telling, explaining and describing at B1.
\end{grammarlens}

\subsection*{Your turn}
\begin{itemize}
\raggedright
  \item \emph{Say it:} \emph{gratiner}
  \item \emph{Say it:} \emph{gratiner}, once more
  \item \emph{Say it:} a sentence of your own with \emph{gratiner}
  \item \emph{Recall:} say the French for the cooking, then the French for to turn golden, then say \emph{gratiner} again
\end{itemize}

\begin{tcolorbox}[breakable,colback=teal!4,colframe=teal!35!black,title={Before you move on}]
What does gratiner mean? (\textquotedblleft{}To brown on top, to cook with a crust\textquotedblright{}.) Say it once more.
\end{tcolorbox}
\section[(dialogue)]{Retour sur les chapitres 324–331}

\label{lesson:FR-R331-revision-324-331}



\begin{quote}
Say the French for to invade and the invasion.
\end{quote}

\subsection*{Your turn}
Say these aloud:

\begin{itemize}
\raggedright
  \item to invade, the invasion, to conquer, the conquest, to defeat
  \item the historian, the ancestor, heritage, the archaeologist, the dig
  \item to free, the liberation, to resist, the resistance, to commemorate
  \item the recipe, the ingredient, the pinch, the spoonful, to preheat
  \item the chopping board, to peel, to chop, to slice thinly, to grate
  \item the salad bowl, to pour, to beat, the whisk, to whisk
  \item to boil, boiling, the cooking pot, to simmer, the casserole dish
  \item the baking tin, to put in the oven, the cooking, to turn golden, to brown on top
\end{itemize}

\begin{tcolorbox}[breakable,colback=teal!4,colframe=teal!35!black,title={Before you move on}]
To invade? (\textbf{envahir}.) The invasion? (\textbf{l'invasion}.) To conquer? (\textbf{conquérir}.) The conquest? (\textbf{la conquête}.) To defeat? (\textbf{vaincre}.) The historian? (\textbf{l'historien}.) The ancestor? (\textbf{l'ancêtre}.) Heritage? (\textbf{le patrimoine}.) The archaeologist? (\textbf{l'archéologue}.) The dig? (\textbf{la fouille}.) To free? (\textbf{libérer}.) The liberation? (\textbf{la libération}.) To resist? (\textbf{résister}.) The resistance? (\textbf{la résistance}.) To commemorate? (\textbf{commémorer}.) The recipe? (\textbf{la recette}.) The ingredient? (\textbf{l'ingrédient}.) The pinch? (\textbf{la pincée}.) The spoonful? (\textbf{la cuillerée}.) To preheat? (\textbf{préchauffer}.) The chopping board? (\textbf{la planche à découper}.) To peel? (\textbf{éplucher}.) To chop? (\textbf{hacher}.) To slice thinly? (\textbf{émincer}.) To grate? (\textbf{râper}.) The salad bowl? (\textbf{le saladier}.) To pour? (\textbf{verser}.) To beat? (\textbf{battre}.) The whisk? (\textbf{le fouet}.) To whisk? (\textbf{fouetter}.) To boil? (\textbf{bouillir}.) Boiling? (\textbf{l'ébullition}.) The cooking pot? (\textbf{la marmite}.) To simmer? (\textbf{mijoter}.) The casserole dish? (\textbf{la cocotte}.) The baking tin? (\textbf{le moule}.) To put in the oven? (\textbf{enfourner}.) The cooking? (\textbf{la cuisson}.) To turn golden? (\textbf{dorer}.) To brown on top? (\textbf{gratiner}.)
\end{tcolorbox}
